\item El método de A. H. Pfund para medir el índice de refracción de un vidrio se ilustra en la figura P34.49. Una de las caras de una placa de grosor $t$ se pinta de blanco, y un pequeño agujero transparente en el punto P sirve como fuente de rayos divergentes cuando la placa se ilumina desde abajo. El rayo PBB' incide con el ángulo crítico y se refleja en su totalidad. En la superficie pintada aparece un círculo oscuro de diámetro $d$ rodeado por una región iluminada (halo).
\begin{enumerate}
    \item Deduzca una ecuación para hallar $n$ en términos de las cantidades medidas $d$ y $t$.
    \item ¿Cuál es el diámetro del círculo oscuro si $n = 1.52$ para una placa de $0.600\text{ cm}$ de espesor?
    \item Si se usa luz blanca, ¿el borde interior del halo brillante está teñido de luz roja o de luz violeta? Explique.
\end{enumerate}